\section{Rによる群間計画の分析}
\label{b20_BetweenDesign_onR}

\subsection{授業内容}
\subsubsection{概要}
前回に続いて，群間計画(Between Design)のデータ分析をRで実行する方法を学ぶ。今回は，二要因群間計画の分散分析(ANOVA)の実行方法と結果の解釈についても学ぶ。Rの\verb|aov()|関数と\verb|lm()|関数でもよいが，\verb|anovakun|関数を使うとより簡単に，詳細な結果を得ることができる。分散分析表から主効果と交互作用の検定結果を読み取り，有意な結果が得られた場合の追加分析(下位検定や単純主効果の検定)の方法についても学ぶ。また，分散分析の結果を適切に可視化する方法についても学ぶ。Rの基本グラフィックス機能やggplot2パッケージを用いた群間の比較を示す棒グラフやボックスプロット，交互作用を示す交互作用プロットの作成方法を理解する。エラーバーの追加方法，軸ラベルやタイトルの設定，凡例の調整など，読みやすく情報量の多いグラフを作成するためのテクニックを学ぶ。実際のデータセットを用いた演習を通じて，群間計画の分析手順と結果の解釈，および可視化方法を身につける。


\subsubsection{コマ主題細目}
\begin{description}
	\item[分散分析用の関数の導入]
	Rで分散分析を行う際に便利な\verb|anovakun|のインストールと基本的な使い方を学ぶ。\verb|anovakun|はCRANに登録されているようなパッケージとは違い，インターネットを経由してソースコードを環境に導入する必要がある。実験デザインの設定の方法についても理解する。

	\item[データの作りかた] 分析に先立って，データセット(ロング型・ワイド型)についての理解が必要である。データファイルは1セルに1情報，重複や集計，オブジェクトや文字列のないデータセットであることが基本であり，一行ごとに1ケースのデータが入り，各列が変数を表す整然データの考え方を学ぶ。可視化のためのロング型と，分散分析の関数に渡すワイド型の変換についても理解する。

	\item[二要因Between計画の実際] 二要因分散分析の実行方法と結果の解釈について学ぶ。\verb|aov()|関数における二要因の公式表記(例：\verb|y ~ factor1 * factor2|)の意味と使い方，交互作用項の指定方法を理解する。分散分析表から主効果と交互作用の検定結果を読み取る方法，有意な交互作用がある場合の解釈と追加分析(単純主効果の検定など)の必要性について学ぶ。また，二要因分散分析の結果を適切に報告する方法についても理解する。実際のデータセットを用いた演習を通じて，二要因群間計画の分析手順と結果の解釈を身につける。
	      \begin{flushright}
		      $\to$ \textcite{kosugi2019}のP.153--160，\textcite{Hashimoto201607}のP.57--90.
	      \end{flushright}

	\item[分析結果の可視化] 分散分析の結果を適切に可視化する方法について学ぶ。Rの基本グラフィックス機能やggplot2パッケージを用いた群間の比較を示す棒グラフやボックスプロット，交互作用を示す交互作用プロットの作成方法を理解する。エラーバーの追加方法，軸ラベルやタイトルの設定，凡例の調整など，読みやすく情報量の多いグラフを作成するためのテクニックを学ぶ。また，可視化することの重要性と，異なるグラフが伝える情報の違いについても理解する。実際のデータセットを用いた演習を通じて，分析結果の効果的な可視化方法を身につける。
	      \begin{flushright}
		      $\to$ \textcite{Kinosady2021}のP.126--138，\textcite{kosugi2019}のP.90--105.
	      \end{flushright}

\end{description}
\subsubsection{キーワード}
\begin{itemize}
	\item \verb|anovakun()|関数の導入と使い方
	\item \verb|anovakun()|関数の各種オプション設定
	\item 交互作用の可視化
	\item ggplot2による結果の表現
\end{itemize}
\subsection{授業情報}
\paragraph{コマの展開方法}
演習
\paragraph{公認心理師標準シラバスにおける位置づけ}
科目番号5 心理学統計法；①心理学で用いられる統計手法；H 推測統計：代表値と散布度をめぐって(2)／科目番号5 心理学統計法；②統計に関する基礎的な知識；E エクセル，R，SPSS入門(3) に相当
\paragraph{心理統計法標準カリキュラム(日本心理学会)における位置づけ}
第11回 分散分析 細目3・4・5，第8回 統計的推定 細目5 に相当
\subsubsection{予習・復習課題}
\paragraph{予習}
前回学んだ群間計画(Between Design)の基本概念と分散分析の理論を復習しておくこと。特に，平方和の分解，F値の意味，主効果と交互作用の概念について理解を確認しておくとよい。また，RとRStudioの基本操作，特にデータの読み込み方法やデータフレームの操作方法について再確認しておくこと。必要に応じて，第3回，第4回の授業で学んだR操作についての資料を見直しておくとよい。

\paragraph{復習}
授業で学んだ内容を定着させるために，以下の課題に取り組むこと：

\begin{enumerate}
\item 同じく提供されたデータセットを用いて，\verb|anovakun|で二要因群間計画の分散分析を実行し，主効果と交互作用の結果を解釈する
\item 分散分析の結果を適切に可視化し，群間の差や交互作用を示すグラフを作成する
\end{enumerate}

これらの課題をRスクリプトファイルにまとめ，期日までに提出すること。時間内に終わらなかった場合は，次回授業までに完成させておくこと。分析結果の適切な解釈と可視化を重視し，統計的に有意な結果だけでなく，効果量や信頼区間についても言及するよう心がけること。不明点があれば，\textcite{kosugi2019}や\textcite{Hashimoto201607}を参照すること。 